\section{Conclusion}
\label{sec:conclusion}

This paper has presented a unified sizing model for user interfaces. The model derives all element geometry from three variables --- intrinsic text lines $n$, wrapping level $w$, and density factor $d$ --- expressed in a base unit $U = \text{fontSize} / 4$.

The key contributions are:

\begin{enumerate}
  \item A closed-form height formula $h = (n \cdot 6 + 2 \cdot d \cdot w) \cdot U$ that covers all height-defined elements.

  \item A four-level wrapping taxonomy that classifies every UI element by boundary semantics, grounded in Gestalt proximity principles~\cite{wertheimer1923,wagemans2012}.

  \item A five-step density scale whose default value reproduces the canonical button heights adopted by MUI, Ant Design, Chakra UI, GitHub Primer, and Adobe Spectrum~\cite{mui,antdesign,chakra,primer,spectrum}.

  \item Em-based spacing emission that makes the model resolution-independent: the same formulas produce correct geometry at any font size without modification.

  \item Context inheritance through the element tree, enabling local overrides that compose without global coordination.

  \item Orthogonality with the tone (color) axis: size, density, and tone operate as three independent parametric models over a shared propagation substrate, allowing arbitrary composition without special-case handling.

  \item A sub-baseline scale for elements below the text-line height, and a dimension table confirming that 69 production patches map to the model without exceptions.
\end{enumerate}

By replacing per-element token tables with parametric expressions, the model ensures that density changes, font-size rescaling, and responsive adaptation propagate uniformly across the entire UI library. The underlying principle is that spatial design should be derived, not declared.
